% Môn: Vật lý
% Lớp: 11
% Chương: Chương IV. Dòng điện. Mạch điện
% Bài: L11C4  Bài 22. Cường độ dòng điện
% Số câu: 24
% App đọc file này trực tiếp — sửa rồi Commit trên GitHub, bấm Đồng bộ GitHub .tex

% L11C4  Bài 22. Cường độ dòng điện

% ===== Câu 1 =====
% ID: L11C4B22-01-TL
% Mức: NB
\dangbt{Chưa có dạng}
\begin{ex}
% ID: L11C4B22-01-TL
% Mức: NB
Quy ước chiều dòng điện là A. chiều dịch chuyển của các electron. B. chiều dịch chuyển của các ion. C. chiều dịch chuyển của các ion âm. D. chiều dịch chuyển của các điện tích dương. $\nguon{ly11 kntt.pdf}$
\choiceTF

\loigiai{
Chọn D.

Quy ước chiều dòng điện là chiều dịch chuyển của các điện tích dương.
}
\end{ex}

% ===== Câu 2 =====
% ID: L11C4B22-01-TN
% Mức: NB
\begin{ex}
% ID: L11C4B22-01-TN
% Mức: NB
Dòng điện trong kim loại là $\nguon{ly11 kntt.pdf}$
\choice
{dòng dịch chuyển của điện tích.}
{\True dòng dịch chuyển có hướng của các điện tích tự do.}
{dòng dịch chuyển có hướng của các hạt mang điện.}
{dòng dịch chuyển có hướng của các ion dương và âm.}
\loigiai{
Chọn B.

Dòng điện trong kim loại là dòng dịch chuyển có hướng của các điện tích tự do.
}
\end{ex}

% ===== Câu 3 =====
% ID: L11C4B22-02-TL
% Mức: NB
\begin{ex}
% ID: L11C4B22-02-TL
% Mức: NB
Khi đặt hiệu điện thế $U = 8\$, $\mathrm{V}$ vào hai đầu một dây dẫn thì dòng điện chạy qua dây có cường độ $I = 0$, $2\,\mathrm{A}$.Nếu tăng hiệu điện thế thêm $4\,\mathrm{V}$ thì dòng điện chạy qua dây dẫn khi đó có cường độ là A. $0,2\,\mathrm{A}$. B. $0,3\,\mathrm{A}$. C. $0,4\,\mathrm{A}$ D. $0,8\,\mathrm{A}\nguon{ly11 kntt.pdf}$
\choiceTF

\loigiai{
Chọn B.

Cường độ dòng điện cần tìm $I_{2} = \dfrac{U_{2}}{R} = \dfrac{U_{2}}{\dfrac{U_{1}}{I_{1}}} = \dfrac{8 + 4}{\dfrac{8}{0,2}} = 0,$ 3 $\mathrm{A}$
}
\end{ex}

% ===== Câu 4 =====
% ID: L11C4B22-02-TN
% Mức: NB
\begin{ex}
% ID: L11C4B22-02-TN
% Mức: NB
Dòng điện không đổi là $\nguon{ly11 kntt.pdf}$
\choice
{dòng điện có chiều không thay đổi theo thời gian.}
{dòng điện có cường độ thay đổi theo thời gian.}
{dòng điện có điện lượng chuyển qua tiết diện thẳng của dây thay đổi theo thời gian.}
{\True dòng điện có chiều và cường độ không thay đổi theo thời gian.}
\loigiai{
Chọn D.

Dòng điện không đổi là dòng điện có chiều và cường độ không thay đổi theo thời gian.
}
\end{ex}

% ===== Câu 5 =====
% ID: L11C4B22-03-TL
% Mức: TH
\begin{ex}
% ID: L11C4B22-03-TL
% Mức: TH
Trong dông sét, một điện tích âm có độ lớn $1\,\mathrm{C}$ được phóng xuống đất trong khoảng thời gian 4 $\cdot$  10^{-4} s. Tính cường độ dòng điện của tia sét đó. $\nguon{ly11 kntt.pdf}$
\choiceTF

\loigiai{
Áp dụng công thức $\text{I} = \dfrac{\text{q}}{\text{t}} = \dfrac{1}{4 \cdot 10^{-4}} = 2500\text{A}$.
}
\end{ex}

% ===== Câu 6 =====
% ID: L11C4B22-03-TN
% Mức: NB
\begin{ex}
% ID: L11C4B22-03-TN
% Mức: NB
Chỉ ra câu sai. $\nguon{ly11 kntt.pdf}$
\choice
{Cường độ dòng điện được đo bằng ampe kế.}
{Để đo cường độ dòng điện, phải mắc nối tiếp ampe kế với mạch điện.}
{Dòng điện chạy qua ampe kê đi vào chốt dương, đi ra chốt âm của ampe kê.}
{\True Dòng điện chạy qua ampe kế đi vào chốt âm, đi ra chốt dương của ampe kế.}
\loigiai{
Chọn D.

Quy ước chiều dòng điện chạy qua ampe kế đi vào chốt dương, đi ra chốt âm của ampe kế.
}
\end{ex}

% ===== Câu 7 =====
% ID: L11C4B22-04-TL
% Mức: TH
\begin{ex}
% ID: L11C4B22-04-TL
% Mức: TH
Trong thời gian 30 giây, có một điện lượng $60\,\mathrm{C}$ chuyển qua tiết diện thẳng của dây dẫn. Tính cường độ dòng điện qua dây và số electron chuyển qua tiết diện thẳng của dây dẫn trong thời gian 2 giây. $\nguon{ly11 kntt.pdf}$
\choiceTF

\loigiai{
- Cường độ dòng điện: $I = \dfrac{\text{\Delta }q}{\text{\Delta }t} = 2\text{A}$.

- Điện lượng chuyển qua tiết diện thẳng của dây dẫn trong thời gian 2 giây:

$\Delta$ q = It = 2.2 = $4\,\mathrm{C}$.

- Số electron chuyển qua tiết diện thẳng của dây dẫn là:

n = [Trong thời gian 30 giây, có một điện lượng $60\,\mathrm{C}$] = 2,5 $\cdot$ 10^{19} electron.
}
\end{ex}

% ===== Câu 8 =====
% ID: L11C4B22-04-TN
% Mức: NB
\begin{ex}
% ID: L11C4B22-04-TN
% Mức: NB
Số electron đi qua tiết diện thẳng của một dây dẫn kim loại trong $1\,\mathrm{s}$, biết điện lượng $30\,\mathrm{C}$ dịch chuyển qua tiết diện đó trong $30\,\mathrm{s}$, là $\nguon{ly11 kntt.pdf}$
\choice
{$3\cdot 10^{18}$.}
{\True $6,25\cdot 10^{18}$.}
{$9\cdot 10^{19}$.}
{$3\cdot 10^{19}$.}
\loigiai{
Chọn B.

Số electron cần tìm $n = \dfrac{30}{1,6\cdot 10^{-19}.30} = 6,25\cdot 10^{18}$
}
\end{ex}

% ===== Câu 9 =====
% ID: L11C4B22-05-TL
% Mức: TH
\begin{ex}
% ID: L11C4B22-05-TL
% Mức: TH
Số electron đi qua tiết diện thẳng của một dây dẫn kim loại trong mỗi giây là $1,25\cdot 10^{19}$. Tính cường độ dòng điện chạy qua dây dẫn và điện lượng đi qua tiết diện đó trong $2$ phút. $\nguon{ly11 kntt.pdf}$
\choiceTF

\loigiai{
Cường độ dòng điện: $\text{I} = \dfrac{\text{\Delta q}}{\text{\Delta t}} = \dfrac{\text{Ne}}{\text{\Delta t}} = \dfrac{1,25 \cdot 10^{19} \cdot 1,6 \cdot 10^{-19}}{1} = 2\text{A.}$ Điện lượng chạy qua tiết diện đó trong 2 phút: q = It = 2.120 = $240\,\mathrm{C}$.
}
\end{ex}

% ===== Câu 10 =====
% ID: L11C4B22-05-TN
% Mức: NB
\begin{ex}
% ID: L11C4B22-05-TN
% Mức: NB
Dòng điện chạy qua một dây dẫn kim loại có cường độ $1\,\mathrm{A}$. Số electron dịch chuyển qua tiết diện thẳng của dây dẫn này trong $2\,\mathrm{s}$ là $\nguon{ly11 kntt.pdf}$
\choice
{$2,5\cdot 10^{19}$.}
{\True $1,25\cdot 10^{19}$.}
{$2\cdot 10^{19}$.}
{$0,5\cdot 10^{19}$.}
\loigiai{
Chọn B.

Số electron cần tìm $n = \dfrac{2.1}{1,6\cdot 10^{-19}} = 1,25\cdot 10^{19}$
}
\end{ex}

% ===== Câu 11 =====
% ID: L11C4B22-06-TL
% Mức: TH
\begin{ex}
% ID: L11C4B22-06-TL
% Mức: TH
Cường độ của dòng điện không đổi chạy qua dây tóc của bóng đèn là $0,64\,\mathrm{A}$. a) Tính điện lượng dịch chuyển qua tiết diện thẳng của dây tóc trong thời gian 2 phút. b) Tính số electron dịch chuyển qua tiết diện thẳng của dây tóc trong khoảng thời gian nói trên. $\nguon{ly11 kntt.pdf}$
\choiceTF

\loigiai{
a) Điện lượng chạy qua tiết diện thẳng của dây tóc trong 2 phút: q = It = 76, $8\,\mathrm{C}$.

b) Số electron chuyển qua tiết diện thẳng của dây tóc là: $\text{n} = \dfrac{\text{q}}{\text{e}} = 4,8 \cdot 10^{20}$ electron.
}
\end{ex}

% ===== Câu 12 =====
% ID: L11C4B22-06-TN
% Mức: NB
\begin{ex}
% ID: L11C4B22-06-TN
% Mức: NB
Trong dây dẫn kim loại có dòng điện không đổi cường độ $2\,\mathrm{mA}$. Trong $1$ phút, số electron chuyển qua một tiết diện thẳng của dây dẫn là $\nguon{ly11 kntt.pdf}$
\choice
{$2\cdot 10^{20}$.}
{$1,2\cdot 10^{19}$.}
{$6\cdot 10^{18}$.}
{\True $7,5\cdot 10^{17}$.}
\loigiai{
Chọn D.

Số lượng electron cần tìm: $n = \dfrac{60.2\cdot 10^{-3}}{1,6\cdot 10^{-19}} = 7,5\cdot 10^{17}$
}
\end{ex}

% ===== Câu 13 =====
% ID: L11C4B22-07-TL
% Mức: TH
\begin{ex}
% ID: L11C4B22-07-TL
% Mức: TH
Tính tốc độ dịch chuyển có hướng của electron trong một dây đồng tiết diện thẳng 1 mm_2 có dòng điện $1\,\mathrm{A}$ chạy qua. Cho biết khối lượng riêng của đồng $\rho = 9$. $103 \mathrm{kg}$ /m_3 và mỗi nguyên tử đồng cho một electron tự do. $\nguon{ly11 kntt.pdf}$
\choiceTF

\loigiai{
Từ khối lượng riêng ta suy ra được số mol trong $1\,\mathrm{m}$ ^{3} đồng là $\dfrac{\rho}{\text{A}}$, với $A$ là nguyên tử lượng của đồng. Số electron tự do trong $1\,\mathrm{m}$ ^{3} đồng là: $\text{n} = \dfrac{\rho}{\text{A}}\text{N}_{\text{A}} = 8,4\cdot 10^{28}\text{electron.}$

Tốc độ của electron tự do là: $v = \dfrac{\text{I}}{\text{Sne}} = 7,4\cdot 10^{-5}\text{m/s}$.
}
\end{ex}

% ===== Câu 14 =====
% ID: L11C4B22-07-TN
% Mức: NB
\begin{ex}
% ID: L11C4B22-07-TN
% Mức: NB
Một dòng điện không đổi chạy qua dây dẫn trong thời gian $10\,\mathrm{s}$, có điện lượng $1,6\,\mathrm{C}$ dịch chuyển qua tiết diện thẳng của dây. Số electron chuyển qua tiết diện thẳng của dây dẫn trong $1\,\mathrm{s}$ là $\nguon{ly11 kntt.pdf}$
\choice
{\True $10^{18}$.}
{$10^{19}$.}
{$10^{20}$.}
{$10^{21}$.}
\loigiai{
Chọn A.

Số electron cần tìm: $n = \dfrac{1,6}{10.1,6\cdot 10^{-19}} = 10^{18}$
}
\end{ex}

% ===== Câu 15 =====
% ID: L11C4B22-08-TN
% Mức: NB
\begin{ex}
% ID: L11C4B22-08-TN
% Mức: NB
Trong thời gian $4\,\mathrm{s}$, điện lượng chuyển qua tiết diện thẳng của dây tóc bóng đèn là $2\,\mathrm{C}$. Cường độ dòng điện qua bóng đèn là $\nguon{ly11 kntt.pdf}$
\choice
{\True $0,5\,\mathrm{A}$.}
{$4\,\mathrm{A}$.}
{$5\,\mathrm{A}$.}
{$0,4\,\mathrm{A}$.}
\loigiai{
Chọn A.

Cường độ dòng điện: $I = \dfrac{q}{t} = \dfrac{2}{4} = 0,$ 5 $\mathrm{A}$
}
\end{ex}

% ===== Câu 16 =====
% ID: L11C4B22-09-TN
% Mức: NB
\begin{ex}
% ID: L11C4B22-09-TN
% Mức: NB
Một dòng điện không đổi, sau 2 phút có một điện lượng $24\,\mathrm{C}$ chuyển qua một tiết diện thẳng của dây dẫn. Cường độ của dòng điện chạy qua dây dẫn là $\nguon{ly11 kntt.pdf}$
\choice
{$1,2\,\mathrm{A}$.}
{$0,12\,\mathrm{A}$.}
{\True $0,2\,\mathrm{A}$.}
{$4,8\,\mathrm{A}$.}
\loigiai{
Chọn C.

Cường độ dòng điện: $I = \dfrac{24}{2.60} = 0,$ 2 $\mathrm{A}$
}
\end{ex}

% ===== Câu 17 =====
% ID: L11C4B22-10-TN
% Mức: NB
\begin{ex}
% ID: L11C4B22-10-TN
% Mức: NB
Một dòng điện không đổi chạy qua dây dẫn có cường độ $2\,\mathrm{A}$ thì sau một khoảng thời gian có một điện lượng $4\,\mathrm{C}$ chuyển qua một tiết diện thẳng của dây dẫn đó. Cùng thời gian đó, với dòng điện $4\,\mathrm{A}$ thì có một điện lượng chuyển qua tiết diện thẳng của dây dẫn đó là $\nguon{ly11 kntt.pdf}$
\choice
{$16\,\mathrm{C}$.}
{$6\,\mathrm{C}$.}
{$32\,\mathrm{C}$.}
{\True $8\,\mathrm{C}$.}
\loigiai{
Chọn D.

Điện lượng cần tìm: $q_{2} = \dfrac{q_{1}}{I_{1}}.I_{2} = \dfrac{4}{2}.4 =$ 8 $\mathrm{C}$
}
\end{ex}

% ===== Câu 18 =====
% ID: L11C4B22-11-TN
% Mức: NB
\begin{ex}
% ID: L11C4B22-11-TN
% Mức: NB
Nếu trong khoảng thời gian $\Delta t = 0$, $1\,\mathrm{s}$ đầu có điện lượng $q = 0$, $5\,\mathrm{C}$ và trong thời gian $\Delta$ t ′ = $0,1\,\mathrm{s}$ tiếp theo có điện lượng q′ = $0,1\,\mathrm{C}$ chuyển qua tiết diện thẳng của vật dẫn thì cường độ dòng điện trong cả hai khoảng thời gian đó là $\nguon{ly11 kntt.pdf}$
\choice
{$6\,\mathrm{A}$.}
{\True $3\,\mathrm{A}$.}
{$4\,\mathrm{A}$.}
{$2\,\mathrm{A}$.}
\loigiai{
Chọn B.

Cường độ dòng điện cần tìm: $I = \dfrac{q_{1} + q_{2}}{t_{1} + t_{2}} = \dfrac{0,5 + 0,1}{0,1 + 0,1} =$ 3 $\mathrm{A}$
}
\end{ex}

% ===== Câu 19 =====
% ID: L11C4B22-12-TN
% Mức: NB
\begin{ex}
% ID: L11C4B22-12-TN
% Mức: NB
Đơn vị của cường độ dòng điện, suất điện động, điện lượng lần lượt là $\nguon{ly11 kntt.pdf}$
\choice
{vôn (V), ampe (A), ampe (A).}
{\True ampe (A), vôn (V), cu lông (C).}
{niutơn (𝑁), fara (F), vôn (V).}
{fara (F), vôn/mét $(\mathrm{V/m})$, jun (J).}
\loigiai{
Chọn B.

Đơn vị của cường độ dòng điện, suất điện động, điện lượng lần lượt là ampe (A), vôn (V), cu lông (C).
}
\end{ex}

% ===== Câu 20 =====
% ID: L11C4B22-13-TN
% Mức: NB
\begin{ex}
% ID: L11C4B22-13-TN
% Mức: NB
Số electron dịch chuyển qua tiết diện thẳng của dây trong thời gian $2\,\mathrm{s}$ là $6,25\cdot 10^{18}$. Khi đó, cường độ dòng điện qua dây dẫn là $\nguon{ly11 kntt.pdf}$
\choice
{$1\,\mathrm{A}$.}
{$2\,\mathrm{A}$.}
{$1,25\,\mathrm{A}$.}
{\True $0,5\,\mathrm{A}$.}
\loigiai{
Chọn D.

Cường độ dòng điện $I = \dfrac{q}{t} = \dfrac{6,25\cdot 10^{18}.1,6\cdot 10^{-19}}{2} = 0,$ 5 $\mathrm{A}$
}
\end{ex}

% ===== Câu 21 =====
% ID: L11C4B22-14-TN
% Mức: NB
\begin{ex}
% ID: L11C4B22-14-TN
% Mức: NB
Số electron đi qua tiết diện thẳng của một dây dẫn kim loại trong $1\,\mathrm{s}$ là $1,25\cdot 10^{19}$. Biết dòng điện không đổi, điện lượng đi qua tiết diện đó trong $15\,\mathrm{s}$ là $\nguon{ly11 kntt.pdf}$
\choice
{$10\,\mathrm{C}$.}
{$20\,\mathrm{C}$.}
{\True $30\,\mathrm{C}$.}
{$40\,\mathrm{C}$.}
\loigiai{
Chọn C.

Điện lượng cần tìm q = 1,25 $\cdot10^{19}.1,6\cdot10^{-19}.15 =30\,\mathrm{C}$
}
\end{ex}

% ===== Câu 22 =====
% ID: L11C4B22-15-TN
% Mức: NB
\begin{ex}
% ID: L11C4B22-15-TN
% Mức: NB
Cường độ dòng điện chạy qua tiết diện thẳng của dây dẫn là 1,5#A.Trong khoảng thời gian $3\,\mathrm{s}$ thì điện lượng chuyển qua tiết diện dây là $\nguon{ly11 kntt.pdf}$
\choice
{$0,5\,\mathrm{C}$.}
{$2,0\,\mathrm{C}$.}
{\True $5,5\,\mathrm{C}$.}
{$5,4\,\mathrm{C}$.}
\loigiai{
Chọn C.

Điện lượng cần tìm q = It = 1,5.3 = 4, $5\,\mathrm{C}$
}
\end{ex}

% ===== Câu 23 =====
% ID: L11C4B22-16-TN
% Mức: NB
\begin{ex}
% ID: L11C4B22-16-TN
% Mức: NB
Cho một dòng điện không đổi chạy qua một dây dẫn, trong $10\,\mathrm{s}$, điện lượng chuyển qua tiết diện thẳng của dây dẫn là $2\,\mathrm{C}$.Sau $60\,\mathrm{s}$ điện lượng chuyển qua tiết diện thẳng đó là $\nguon{ly11 kntt.pdf}$
\choice
{$6\,\mathrm{C}$.}
{\True $12\,\mathrm{C}$.}
{$60\,\mathrm{C}$.}
{$20\,\mathrm{C}$.}
\loigiai{
Chọn B.

Điện lượng cần tìm $q = \dfrac{2}{10}.60 =$ 12 $\mathrm{C}$
}
\end{ex}

% ===== Câu 24 =====
% ID: L11C4B22-17-TN
% Mức: TH
\begin{ex}
% ID: L11C4B22-17-TN
% Mức: TH
Cường độ dòng điện được xác định theo biểu thức nào sau đây? $\Delta$ q $\Delta$ t $\Delta$ q $\nguon{ly11 kntt.pdf}$
\choice
{$I = \Delta q$. $\Delta$ t.}
{\True $I = \Delta t$.}
{$I = \Delta q$.}
{$I = e$.}
\loigiai{
Chọn B.

Cường độ dòng điện được xác định theo biểu thức $I = \dfrac{\text{\Delta }q}{\text{\Delta t}}$.
}
\end{ex}
